\documentclass[12pt]{article}
\usepackage[T1]{fontenc}
\usepackage[utf8]{inputenc}
\usepackage{lmodern}
\usepackage{textcomp}
\usepackage{enumitem}
\usepackage{geometry}
\geometry{margin=1in}
\begin{document}
\begin{center}
Answer Key - Constitutional Law - Graduate Level Assessment 10 \
\small (Answers are ordered exactly as in the question paper.)
\end{center}

\section*{Multiple Choice}
\begin{enumerate}[label=\arabic*.]
\item \textbf{Answer:} D
\\ \textit{Options:} A) It is jurisdiction based on the country where the legal person was Registered; B) It is jurisdiction based on the nationality of the offender; C) It is jurisdiction based on the relationship between the victim and the offender; D) It is jurisdiction based on the nationality of the victims; E) It is jurisdiction based on the severity of the offence
\\ \textit{Note:} Passive personality jurisdiction refers to the legal principle that allows a state to exercise jurisdiction over a crime committed against its nationals, regardless of where the crime occurred, thus prioritizing the victims' nationality as the basis for legal authority.
\item \textbf{Answer:} A
\\ \textit{Options:} A) There is no universal agreement on what constitutes natural law.; B) Natural law is an arbitrary concept.; C) Natural law does not apply to all cultures and societies.; D) Natural law is backward-looking.; E) There is no such thing as a social contract.
\\ \textit{Note:} Hume's critique of natural law centers on the observation that varying cultural, moral, and philosophical perspectives lead to a lack of consensus on what principles should be universally accepted as natural law, undermining its validity as an objective standard.
\item \textbf{Answer:} A
\\ \textit{Options:} A) The UK Constitution is a direct copy of the US Constitution; B) The UK Constitution is codified and contained in a single document; C) The UK Constitution can be changed by a simple majority vote in Parliament; D) The UK Constitution gives the judiciary the power to overturn acts of parliament; E) The UK Constitution is based on a Bill of Rights
\\ \textit{Note:} The statement is incorrect because the UK Constitution is not a direct copy of the US Constitution; rather, it is an uncodified and flexible constitution based on statutes, common law, and conventions, reflecting the UK's unique legal and political history.
\item \textbf{Answer:} C
\\ \textit{Options:} A) The end of the Cold War; B) The Bolshevik Revolution; C) The international recognition of human rights after World War II; D) The invention of the Internet; E) The rise of Fascism
\\ \textit{Note:} The international recognition of human rights after World War II reinvigorated natural law by emphasizing inherent human dignity and moral principles that transcend positive law, aligning legal frameworks with universal ethical standards.
\item \textbf{Answer:} A
\\ \textit{Options:} A) Socialism.; B) Legitimacy.; C) Authority.; D) Democracy.; E) Bureaucracy.
\\ \textit{Note:} Weber's analysis of formally rational law emphasizes the role of socialism in promoting a structured legal system that prioritizes efficiency and predictability, reflecting how economic and social structures influence legal development.
\end{enumerate}

\section*{True / False}
\begin{enumerate}[label=\arabic*.]
\setcounter{enumi}{5}
\item \textbf{Answer:} False
\\ \textit{Options:} A) True; B) False
\\ \textit{Note:} The statement is false because while regulating the use of public spaces can contribute to social order, it is primarily aimed at maintaining public safety and order rather than directly ensuring individual freedom and social justice, which can sometimes be at odds with such regulations.
\item \textbf{Answer:} True
\\ \textit{Options:} A) True; B) False
\\ \textit{Note:} Under international law, all states have jurisdiction over crimes committed on board vessels, as they can assert jurisdiction based on the flag state principle, which allows the state under whose flag the vessel is registered to exercise legal authority, as well as the principle of universal jurisdiction for certain serious offenses.
\item \textbf{Answer:} True
\\ \textit{Options:} A) True; B) False
\\ \textit{Note:} Parfit argues that by addressing the needs of the poor through equitable resource distribution and social policies, society can reduce disparities and promote greater overall equality.
\item \textbf{Answer:} False
\\ \textit{Options:} A) True; B) False
\\ \textit{Note:} The definition of law as a set of principles recognized by the State was first proposed by legal positivists before Hart, notably by thinkers such as Jeremy Bentham, making the statement false.
\item \textbf{Answer:} True
\\ \textit{Options:} A) True; B) False
\\ \textit{Note:} Cognitivism asserts that ethical truths are objective and universal, existing independently of personal beliefs, which directly contrasts with ethical relativism that argues ethical standards, including human rights, are shaped by cultural and individual perspectives.
\end{enumerate}

\section*{Long Form Response}
\begin{enumerate}[label=\arabic*.]
\setcounter{enumi}{10}
\item \textbf{Answer:} The concept of "excusable neglect" refers to a legal standard allowing for the late filing of documents, such as a notice of appeal, when a party can demonstrate that their failure to file on time was due to circumstances beyond their control. In the case of a litigant who caused a car crash by running a red light and subsequently faced hospitalization for two months, this serious personal circumstance could justify a delay in filing an appeal. The hospitalization could be seen as a compelling reason for the litigant's inability to timely respond to the court's judgment, as it directly impacted their capacity to engage in legal proceedings. Courts typically evaluate whether the neglect was reasonable under the circumstances, and significant health crises, particularly those resulting from accidents, may warrant leniency in acknowledging the litigant's inability to act promptly. Thus, the combination of the accident and subsequent recovery period could substantiate a claim of excusable neglect in the context of the appeal process.
\\ \textit{Note:} correct because it identifies "excusable neglect" as a legal standard that permits late filings due to unforeseen and serious circumstances, such as hospitalization from an accident, which can clearly impede a person's ability to comply with filing deadlines.
\item \textbf{Answer:} The enforceability of sales offers, particularly in the context of a merchant's promise to keep an offer open, is primarily governed by the Uniform Commercial Code (UCC). When a merchant, such as the one in the scenario, emails a buyer to offer a widget for \$35,000 with a commitment to keep the offer open for ten days, this constitutes an offer that may create a binding obligation if it meets certain conditions. Under UCC § 2-205, a merchant's signed offer to sell goods, which includes a promise to keep the offer open, cannot be revoked for the specified time, provided the offer is in writing. However, if the offer lacks consideration, it may be subject to revocation prior to acceptance. Thus, while the ten-day promise enhances the offer's enforceability, the absence of consideration or a written agreement can ultimately allow for its revocation.
\\ \textit{Note:} The answer correctly identifies that under UCC § 2-205, a merchant's signed offer to keep an offer open for a specified period can create a binding obligation, thus highlighting the enforceability of sales offers and the conditions that govern their revocation.
\end{enumerate}

\vfill
\noindent\textit{End of Answer Key}
\end{document}